\subsubsection{Arquitectura general de la red privada}

La red privada \gls{5g} del \gls{iteamupv} constituye la infraestructura de comunicaciones sobre la que se soporta la capa robótica descrita en la Sección \ref{sec:capa_robotica}. Esta red presenta una arquitectura completa en modo \gls{5gsa}, con despliegue \textit{indoor} y \textit{outdoor}, múltiples bandas de frecuencia y diversas implementaciones de núcleo \gls{5g} \cite{ITEAM2025Private5G}.

Desde el punto de vista funcional, la arquitectura se estructura en cuatro bloques principales:

\begin{enumerate}
    \item \textbf{Radio Access Network (RAN)}

    La \gls{ran} es el subsistema encargado de proporcionar conectividad radio entre los \gls{ues} y el núcleo \gls{5g}. Entre las características más relevantes del equipamiento empleado destacan:

    \begin{itemize}
        \item Ancho de banda de hasta 100 MHz.
        \item Configuración flexible en \textit{Time Division Duplex} (TDD).
        \item Soporte MIMO 4x4.
        \item Arquitectura funcionalmente separada \gls{cu_du_ru}.
    \end{itemize}

    La separación funcional \gls{cu_du_ru} permite desacoplar el procesamiento de la capa física de las capas superiores, reducir la latencia situando la \gls{du} próxima a la unidad radio y virtualizar la \gls{cu} sobre infraestructura de propósito general. Asimismo, el soporte de \textit{radio slicing} mediante asignación dinámica de recursos físicos (PRBs) permite priorizar determinados flujos de tráfico en función de los requisitos de cada aplicación.

    En aplicaciones de robótica móvil, la \gls{ran} resulta crítica al garantizar un \textit{uplink} robusto para la transmisión de datos sensoriales y minimizar el \textit{jitter} en escenarios de teleoperación.

    \item \textbf{Núcleo \gls{5g} (\gls{5g} Core)}

    El núcleo \gls{5g} implementa las funciones del plano de control (\textit{Control Plane}) y del plano de usuario (\textit{User Plane}). El núcleo sigue una arquitectura \gls{sba}, propia de \gls{5gsa}, lo que permite la virtualización y desacoplamiento de funciones de red sobre infraestructura de propósito general.

    Estas funciones pueden ejecutarse de forma virtualizada tanto en infraestructura \textit{Edge} como en \textit{Cloud}, tal y como se representa en la Figura~\ref{fig:arquitectura_5Gsa}.

    La posibilidad de instanciar múltiples \gls{upf} habilita mecanismos de \textit{network slicing} a nivel de \textit{core}, permitiendo:

    \begin{itemize}
        \item Aislar tráfico por aplicación.
        \item Implementar políticas de QoS diferenciadas.
        \item Aproximar el plano de usuario al \textit{Edge}.
    \end{itemize}

    En sistemas distribuidos basados en \gls{ros} con estrategias de \textit{offloading}, la proximidad de la \gls{upf} a los nodos \textit{Edge} reduce la latencia \gls{e2e} en la transmisión de \textit{topics} de alta tasa, lo que resulta determinante en aplicaciones sensibles al retardo.

    \item \textbf{Infraestructura Edge asociada al plano de usuario}

    La infraestructura \textit{Edge} constituye el punto intermedio del \textit{continuum} 	\textit{Robot--Edge--Cloud} \cite{CrespoAguado2024IoTEdgeCloud}. Asociada al plano de usuario, permite ejecutar tareas de procesamiento con requisitos estrictos de latencia.

    La infraestructura \textit{edge} puede alojar instancias virtualizadas de la \gls{upf}, aproximando el plano de usuario a los dispositivos finales y reduciendo la latencia \gls{e2e}. Esta integración refuerza el carácter distribuido y programable de la arquitectura, alineado con el paradigma de \gls{nfv}.

    En el contexto de robótica móvil, el \textit{edge} posibilita:

    \begin{itemize}
        \item Teleoperación inmersiva.
        \item Procesamiento en tiempo real de vídeo.
        \item \textit{Offloading} de algoritmos de \gls{slam} de elevada complejidad computacional.
        \item Inferencia de modelos de inteligencia artificial.
    \end{itemize}

    Este enfoque reduce la carga computacional del robot y evita el envío continuo de grandes volúmenes de datos al \textit{cloud}.

    \item \textbf{Interconexión con cloud}

    La interconexión con infraestructuras \textit{cloud} permite ejecutar tareas computacionalmente intensivas no críticas en latencia, tales como:

    \begin{itemize}
        \item Procesamiento Big Data.
        \item Entrenamiento de modelos de aprendizaje automático.
        \item Almacenamiento histórico de datos.
    \end{itemize}

    Este diseño desacopla los procesos sensibles a la latencia (gestionados en \textit{edge}) de aquellos con mayores requerimientos computacionales pero menor criticidad temporal (gestionados por el \textit{cloud}) \cite{ITEAM2025Private5G}.
\end{enumerate}

Como se muestra en la Figura~\ref{fig:arquitectura_5Gsa}, la arquitectura integra la red privada \gls{5g} con la infraestructura \textit{edge} y la interconexión \textit{cloud} bajo un esquema de orquestación dual (\textit{Network Orchestration} y \textit{Edge Orchestration}). Esta separación permite la gestión dinámica de recursos, la instanciación de funciones de red virtualizadas y el despliegue automatizado de aplicaciones distribuidas.

\begin{figure}[H]
    \centering
    \includegraphics[width=0.85\linewidth]{figures/arquitectura_5Gsa.png}
    \caption{Arquitectura funcional de la red privada \gls{5gsa} integrada con infraestructura \textit{Edge} y \textit{Cloud}, donde se muestra la virtualización de funciones de red y aplicaciones, así como los mecanismos de orquestación \cite{ITEAM2025Private5G}.}
    \label{fig:arquitectura_5Gsa}
\end{figure}

La adopción de una red privada \gls{5g} como capa de red del sistema experimental responde a la necesidad de garantizar aislamiento lógico, control determinista de \gls{qos} y segmentación mediante \textit{network slicing}, capacidades que no están disponibles de forma nativa en tecnologías inalámbricas convencionales como IEEE 802.11 (WiFi), especialmente en entornos con tráfico heterogéneo y requisitos de baja latencia.

Desde una perspectiva sistémica, la red privada no actúa únicamente como medio de transporte de datos, sino como un componente activo del sistema distribuido, alineado con el paradigma 	\textit{Robot--Edge--Cloud} y con los requisitos de aplicaciones avanzadas como Digital Twins y teleoperación inmersiva.

% Finalmente, la red privada opera en las bandas n40 y n78, lo que permite evaluar compromisos entre cobertura, capacidad y latencia.